% ============================================
% MEM-06
% Gerado de molambudos.md — build_miolo.py
% ============================================
\fragtitle{mem}{M}{MEM}{06}{O Curral do Governo}
\markboth{MEM-06 --- O Curral do Governo}{}

\noindent \textit{Os Molambudos --- Campos de Concentração do Ceará}

\noindent \textbf{1915 --- Campos de Concentração do Ceará}

\vspace{0.5em}\noindent\rule{\textwidth}{0.5pt}\vspace{0.5em}

\noindent O governo chamava de \textquotedblleft{}currais\textquotedblright{}.

\noindent A palavra certa. A gente era gado.

\noindent O Campo do Alagadiço, nos arredores de Fortaleza, era uma área cercada de arame farpado, com uns duzentos metros de largura por talvez trezentos de comprimento. Não tinha cobertura. Não tinha parede. Não tinha nada além de terra batida, sol, e o cheiro de gente acumulada.

\noindent No primeiro dia, eu contei as pessoas. Cheguei a trezentas e poucas antes de perder a conta. Depois chegou mais gente. Depois mais. No fim da primeira semana, a gente se apertava como bicho no transporte.

\noindent Todo dia chegavam soldados com mais retirantes. Eles vinham de toda parte: Quixeramobim, Boa Viagem, Senador Pompeu, Quixadá, Canindé. Os soldados chamavam a gente de \textquotedblleft{}molambudos\textquotedblright{}. A palavra saía da boca deles como se fosse o nome de uma praga.

\noindent \textit{\textquotedblleft{}Lá vem mais molambudos.\textquotedblright{}}

\noindent \textit{\textquotedblleft{}Molambudo, sai daí.\textquotedblright{}}

\noindent \textit{\textquotedblleft{}Dá comida pra esses molambudos antes que morram tudo.\textquotedblright{}}

\vspace{0.5em}\noindent\rule{\textwidth}{0.5pt}\vspace{0.5em}

\noindent A comida era uma vez por dia.

\noindent Um caldo ralo. Os soldados chamavam de \textquotedblleft{}sopa do governo\textquotedblright{}. Era água com farinha e algumas folhas boiando. Eu tomava a minha e ainda lambia a cuia. Minha mãe dava a dela pra mim. Eu dizia que não. Ela dizia que já tinha comido. Mentira. Eu via a tosse sacudir o corpo vazio dela.

\noindent O governo mandava boi de vez em quando. Não para alimentar a gente --- para alimentar os soldados. A gente ficava olhando eles assarem a carne. O cheiro passava por cima da cerca e entrava no nosso estômago. Alguns choravam. Outros batiam no arame farpado. Os soldados riam.

\noindent \textit{\textquotedblleft{}Deixa eles cheirar\textquotedblright{}\textit{, um deles disse. }\textquotedblleft{}Fome não se cura com cheiro.\textquotedblright{}}

\noindent Ele tinha razão. Fome não se cura com cheiro. Nem com choro. Nem com oração. A fome só se cura com comida.

\noindent E não tinha comida.

\vspace{0.5em}\noindent\rule{\textwidth}{0.5pt}\vspace{0.5em}

\noindent As mortes começaram na segunda semana.

\noindent Primeiro os velhos. Depois as crianças. Depois os adultos que já estavam doentes antes de chegar.

\noindent Os soldados tinham uma rotina. De manhã, abriam a cerca, entravam, andavam entre a gente, e apontavam os que tinham morrido durante a noite. Carregavam os corpos pelos braços e pernas --- um de cada lado --- e levavam para a vala.

\noindent A vala ficava a duzentos metros. A gente via os urubus.

\noindent Minha mãe tossia.

\vspace{0.5em}\noindent\rule{\textwidth}{0.5pt}\vspace{0.5em}

\noindent Uma noite, eu sonhei com comida.

\noindent Era um sonho bom. Eu estava na casa da minha avó, que tinha morrido antes de eu nascer, mas no sonho ela estava viva. Ela fazia cuscuz. O cuscuz saía da panela e subia até o céu. Eu abria a boca e ele caía dentro de mim como chuva.

\noindent Acordei com fome.

\noindent Uma fome que doía. Não no estômago. No corpo inteiro. Os ossos doíam. A boca doía. Os dentes doíam. Eu olhei para as mãos e vi os ossos por baixo da pele.

\noindent \textit{\textquotedblleft{}Mãe\textquotedblright{}\textit{, eu disse. }\textquotedblleft{}Eu quero ir embora.\textquotedblright{}}

\noindent Ela não respondeu.

\noindent \textit{\textquotedblleft{}Mãe?\textquotedblright{}}

\noindent Ela estava deitada de lado. O corpo dela estava quente. A respiração era curta. Rápida. Como a de um passarinho.

\noindent \textit{\textquotedblleft{}Vai\textquotedblright{}\textit{, ela disse. }\textquotedblleft{}Vai, Joaquim. Vai.\textquotedblright{}}

\noindent \textit{\textquotedblleft{}Pra onde?\textquotedblright{}}

\noindent Ela tossiu. A tosse sacudiu o corpo inteiro. Quando a tosse passou, os lábios dela estavam vermelhos.

\noindent \textit{\textquotedblleft{}Pra longe daqui.\textquotedblright{}}

\noindent \textit{\textquotedblleft{}Eu não sei pra onde ir.\textquotedblright{}}

\noindent Ela fechou os olhos.

\noindent \textit{\textquotedblleft{}Tu sabe.\textquotedblright{}}

\noindent E eu sabia.

\noindent Eu não queria saber. Mas sabia.

\noindent A vala. A única coisa que ficava além do arame farpado. A única direção que não exigia caminhar. A vala. Lá tinha carne.

\noindent Eu não fui naquela noite.

\noindent Fui na seguinte.\NE{A existência do Campo do Alagadiço, próximo a Fortaleza, durante a seca de 1915 é documentada. Joaquim, sua família e os episódios sensoriais desta passagem são ficcionais; detalhes de escala, rotina e vala pertencem à reconstituição narrativa e não equivalem a testemunho arquivístico.}

\vspace{0.5em}\noindent\rule{\textwidth}{0.5pt}\vspace{0.5em}

\noindent O Campo do Alagadiço foi criado durante a seca de 1915, nos arredores de Fortaleza. Registros históricos indicam que milhares de retirantes foram confinados ali, mas as fontes consultadas não permitem fixar, para esta narrativa, um total preciso de internados ou mortos.

\noindent Eu só sabia que muitos não saíam.

\noindent Eu era um dos que sobraram.

\noindent As vezes eu penso: \textit{por que eu?}

\noindent Aí a criatura ri dentro de mim.

\noindent E eu lembro que eu não sobrevivi \textit{apesar} da vala.

\noindent Eu sobrevivi \textit{por causa} dela.

\vspace{0.5em}
\noindent\textit{\textbf{↪ Rotas:} \rota{MEM-07} \rotasep \rota{MEM-01}}
